% =============================================================================
\section{Introduction}
% =============================================================================

Everyday users increasingly interact with AI assistants for writing,
summarization, planning, and information search.
Two broad classes of model are now practically available: \textit{large language
models} (LLMs), accessed via cloud APIs and characterized by high accuracy and
robustness, and \textit{small language models} (SLMs), which run locally on
commodity hardware, offer stronger privacy guarantees, and operate without
internet connectivity.
The pragmatic question---\textit{for a given daily task, is an SLM good enough?}
---has real stakes for sustainable computing and equitable access, yet it has no
clean empirical answer.

The difficulty is not purely technical.
Observed differences in output quality between an SLM and an LLM are
\textit{entangled} with at least two confounds: (i)~\textit{prompting-skill
variance}, because users write different prompts for the two systems and SLMs are
more brittle to underspecified input~\cite{zamfirescu2023johnny,
subramonyam2024gulf}; and (ii)~\textit{interface differences}, because most
SLM deployments offer a bare chat box while commercial LLM products embed
suggestion, feedback, and correction affordances.
Letting users ``freestyle'' prompts in different tools and comparing outputs
conflates model capability with user skill and UI support---making it impossible
to attribute gaps to the model.

This paper addresses both confounds through a single shared interface,
\textit{Easy\_SLM}, that routes the same user specification to either an SLM
(via effGen~\cite{effgen}) or a cloud LLM (OpenAI / Anthropic), and that
provides \textit{light scaffolding} to control Gulf-of-Envisioning
variance~\cite{subramonyam2024gulf} in both conditions.
Our work is motivated by three concerns.
\textbf{Sustainability:} ``Use a smaller model'' is not actionable without
knowing when an SLM is genuinely sufficient~\cite{crawford2024rematerializing}.
\textbf{Resilient access:} Cloud dependency creates brittleness; the 2024
CrowdStrike outage disrupted AI-assisted care in hundreds of hospitals, and
offline-capable SLMs offer essential fallback~\cite{tully2025crowdstrike}.
\textbf{Democratization:} Novice users face a Gulf of Envisioning with SLMs
because these models are less forgiving of vague intent; scaffolding shifts
``prompting as skill'' to ``prompting as structure''~\cite{zamfirescu2023johnny}.

We address four research questions:
\textbf{RQ1} (Capability gap): How does SLM performance compare to LLM
performance on daily tasks under objective metrics (ROUGE, F1, rule-based
pass rates)?
\textbf{RQ2} (``Good enough'' boundary): For which tasks can SLM outputs be
considered ``good enough'' based on combined objective scores and user
preference?
\textbf{RQ3} (Prompting confound): To what extent does a standardized
prompting interface reduce prompting-skill variance in SLM--LLM comparisons?
\textbf{RQ4} (Affordances): How do instruction decomposition and
constraint-checking affordances affect output quality, format violations, and
perceived workload across conditions?

Our contributions are: (1)~a formative Reddit study identifying 19 everyday
tasks shared across SLM and LLM user communities; (2)~the Easy\_SLM system
with four design goals operationalizing fair paired comparison; and (3)~user
study evidence on capability gaps, ``good enough'' task boundaries, and
scaffold effectiveness.
